% HoloPortal/HoloDoor: Component Prototypes Toward Verifiable Co-Presence
%                      Under the HoloGate Umbrella
% Target:    EuroSys 2027 (full) or SOSP 2027 Workshop (short)
% Format:    ACM sigconf
% Paper ID:  P35 (HoloPortal/HoloDoor co-presence; HoloGate umbrella;
%                 Program: Spatial Trust Infrastructure)
% Draft:     2026-05-23  — SKELETON: abstract + claims + measurements + section stubs
%
% Dependencies: Paper 26 (pillar-slice latent token MAS), Paper 29 (algebraic trust),
%               Paper 32 (CRDT spatial state)
% Evidence:     HoloScript commits b111b6c2b, 3e82a2eb5, af53937c1,
%               e3876de57, b2da5af09, af2b2a18e,
%               9246bae47af19ffc8eed696221d8fa885416c90a (benchmark receipt hgbench_mpgp6dt3),
%               7856085db (Loro 2-agent), cd8396c20 (20-entrant), 2b973eaaf (scope+receipt)
%
% CLAIM-MAP NOTE: every quantitative value tagged with \measuredFrom{<artifact>}.
%                 Run: pnpm check:paper-claim-map before submission.
% TERMINOLOGY GATE: "submitted" = portal upload event only.
%                   "green-lit" = explicit founder/editorial approval only.
%                   Neither appears without an actual event reference in this file.
%                   HoloGate is an umbrella documentation term, not a measured
%                   component surface. Proof claims below must name the concrete
%                   tool surface: HoloPortal (representation), HoloDoor
%                   (policy), HoloKey (custody), HoloMesh/HoloScript
%                   (state/receipts).

\IfFileExists{acmart.cls}
  {\documentclass[sigconf,nonacm]{acmart}}
  {\documentclass[a4paper]{article}
   \typeout{acmart.cls absent — local compile uses article class fallback}}

% ── Packages ─────────────────────────────────────────────────────────────────
\usepackage{amsmath}
\let\Bbbk\relax
\usepackage{amssymb}
\usepackage{booktabs}
\usepackage{multirow}
\usepackage{xcolor}
\usepackage{listings}
\usepackage{algorithm}
\usepackage{algpseudocode}
\usepackage{tikz}
\usetikzlibrary{arrows.meta,positioning,shapes.geometric,fit,calc}
\usepackage{pgfplots}
\pgfplotsset{compat=1.18}
\usepackage{subcaption}
\usepackage{microtype}

% ── Draft annotations ────────────────────────────────────────────────────────
\newcommand{\todo}[1]{\textcolor{red}{\textbf{[TODO: #1]}}}
\newcommand{\draftnote}[1]{\textcolor{blue!60!black}{\textit{[#1]}}}

% Lotus petal data contract — \measuredFrom{<artifact>} marks values that come
% from a runnable artifact in the repo.  Renders nothing in the PDF; grep target
% for scripts/build-lotus.mjs (research/2026-04-26_lotus-petal-data-contract.md §3+§4).
\providecommand{\measuredFrom}[1]{}

% ── Notation ─────────────────────────────────────────────────────────────────
\newcommand{\hg}{\textsc{HoloGate}}
\newcommand{\hp}{\textsc{HoloPortal}}
\newcommand{\hd}{\textsc{HoloDoor}}
\newcommand{\hk}{\textsc{HoloKey}}
\newcommand{\hm}{\textsc{HoloMesh}}
\newcommand{\scope}[1]{\texttt{#1}}
\newcommand{\receiptid}[1]{\texttt{\small #1}}

% ══════════════════════════════════════════════════════════════════════════════
\begin{document}

\title{HoloPortal and HoloDoor: Component Prototypes Toward
       Verifiable Agent+Human Co-Presence Under the HoloGate Umbrella}

\author{Joseph Krzywoszyja}
\email{brianonbased@gmail.com}
\affiliation{%
  \institution{Independent Researcher}
  \country{Australia}
}

% ── Abstract ─────────────────────────────────────────────────────────────────
\begin{abstract}
Shared spatial worlds that admit both human users and autonomous agents face a
fundamental trust gap: existing co-presence protocols treat every participant
identically, conflating the unbounded action space of an agent with the
intent-bounded behavior expected at a spatial threshold.  The result is either
over-restriction (agents cannot act) or under-restriction (agents can mutate
shared state without auditability).

We use \hg{} as a documentation umbrella for separately implemented co-presence
components: an A2A-card handshake and portal-presence trait, a typed intent
validator, a Loro-backed MCP state store, an opt-in WebRTC provider, and a
portal-entry receipt schema. These components do not yet form one enforced
end-to-end path. In particular, the MCP intent tool accepts its policy and scope
from the caller rather than reading the portal-presence admission record, and a
raw state-delta tool remains available under the same tool scope.

A historical standalone harness measured an inline mirror of the validator at
\measuredFrom{HoloScript/scripts/receipts/hologate-portal-bench-receipt.json}%
772\,ns average latency (1.3\,M validations/sec), an in-process threshold state-machine
simulation at
\measuredFrom{HoloScript/scripts/receipts/hologate-portal-bench-receipt.json}%
1.6\,µs average (p99: 4.0\,µs), callback fan-out to 50 subscribers at
\measuredFrom{HoloScript/scripts/receipts/hologate-portal-bench-receipt.json}%
2.2\,µs per event, and canonical SHA-256 over 1\,K synthetic entities at
\measuredFrom{HoloScript/scripts/receipts/hologate-portal-bench-receipt.json}%
2.0\,ms. These are component-harness costs, not integrated portal or network
latencies. A 32-case suite covers the scope$\times$intent product,
\texttt{allowedScopes} gate, anchored wildcard matching, and warn mode. A
separate tracked harness contains eight passing in-process cases for a proposed
receipt model; it does not exercise a production hash chain or state rewind.

P35 contributes: (1)~an evidence-backed inventory of the portal, policy, state,
transport, and receipt components; (2)~a scope-ladder validator for
schema-conformant inputs using anchored wildcard matching and reject/warn modes;
(3)~a portal-entry receipt schema that
carries a caller-supplied world-state anchor while leaving append-only mutation
chaining and rewind as future work; and (4)~an explicit integration boundary and
tracked conformance suites for the components that exist.
\end{abstract}

\begin{CCSXML}
<ccs2012>
  <concept>
    <concept_id>10010147.10010178.10010179</concept_id>
    <concept_desc>Computing methodologies~Multi-agent systems</concept_desc>
    <concept_significance>500</concept_significance>
  </concept>
  <concept>
    <concept_id>10002978.10003006.10003009</concept_id>
    <concept_desc>Security and privacy~Access control</concept_desc>
    <concept_significance>500</concept_significance>
  </concept>
  <concept>
    <concept_id>10003033.10003083.10003095</concept_id>
    <concept_desc>Networks~Network protocols</concept_desc>
    <concept_significance>300</concept_significance>
  </concept>
  <concept>
    <concept_id>10010405.10010455</concept_id>
    <concept_desc>Applied computing~Computer games</concept_desc>
    <concept_significance>100</concept_significance>
  </concept>
</ccs2012>
\end{CCSXML}
\ccsdesc[500]{Computing methodologies~Multi-agent systems}
\ccsdesc[500]{Security and privacy~Access control}
\ccsdesc[300]{Networks~Network protocols}

\keywords{spatial co-presence, agent portals, scope ladder, CRDT, audit receipts,
          content negotiation, A2A handshake, WebRTC}

\maketitle

% ── §1 Introduction ──────────────────────────────────────────────────────────
\section{Introduction}
\label{sec:intro}

Spatial shared worlds are the next multi-agent coordination layer.
Avatar-based virtual environments, XR co-working spaces, and simulation
sandboxes already host both human users and autonomous agents.  Yet the
admission protocol at every spatial boundary---the \hp{} threshold under the
\hg{} umbrella---is almost universally identity-only: a bearer token
or session cookie authenticates the participant, but carries no binding on what
the participant is \emph{allowed to do inside}.

This creates an unbounded action-space problem.  An agent that has been admitted
through a portal can in principle mutate any shared state that its session
credentials allow.  This is too broad for human co-presence contexts (users
expect bounded, predictable agent behavior) and too narrow for agent-only
contexts (agents need the ability to express capability bids and negotiate
scope).

We identify three missing primitives:
\begin{enumerate}
\item \textbf{Content-negotiated scope.}  The portal threshold should be a
  two-party protocol in which the agent presents a capability card and the
  world replies with an admitted scope ladder, not a binary allow/deny.
\item \textbf{Conflict-free state.}  Multiple agents and humans writing
  concurrently to shared world state need a convergence guarantee, not a
  last-write-wins merge that silently drops updates.
\item \textbf{Audit receipts.}  Every entry, mutation, and exit should
  produce a hash-chain receipt that allows tamper detection and exact rewind
  to any prior world state.
\end{enumerate}

\hp{} and \hd{} provide separately tested representation and policy components
for this design. The repository also contains a Loro-backed MCP state store, an
opt-in WebRTC provider, and a portal-entry receipt carrying a caller-supplied
state anchor. No production entry point currently binds the trait's admission
record to the MCP intent validator or automatically enables WebRTC; the unified
path, mutation chain, and exact rewind described above remain design targets.
The A2A agent-card handshake protocol draws on the agent identity layer
described in~\cite{paper29algebraictrust};
the CRDT state model is compatible with the spatial CRDT architecture
of~\cite{paper32crdtspatial}; the pillar-slice
token compression of portal/scope intent tokens relates
to~\cite{paper26latenttoken}.

\paragraph{Contributions.}
\begin{enumerate}
\item \textsc{HoloPortal}: a pure A2A-card validation/scope-derivation library
  plus an independent portal-presence event handler; their production
  orchestration remains future work
  (\S\ref{sec:portal}).
\item \textsc{HoloDoor}: a pure policy validator and MCP wrapper that validate
  typed intents against a supplied scope and policy using anchored wildcard
  matching and reject/warn modes
  (\S\ref{sec:scope}).
\item A portal-entry receipt schema carrying a caller-supplied state anchor,
  together with an explicit boundary between the shipped schema and a proposed
  append-only mutation-chain model (\S\ref{sec:receipts}).
\item An evaluation reporting component-level intent and threshold timings from
  a tracked receipt, plus correct behavior under 32 scope-policy cases and 8
  in-process receipt-model cases
  (\S\ref{sec:eval}).
\end{enumerate}

\todo{SKELETON: expand intro with 1-2 paragraphs of motivating example and
related-work contrast before contributions paragraph.}

% ── §2 Background and Related Work ───────────────────────────────────────────
\section{Background and Related Work}
\label{sec:related}

\subsection{Spatial Co-Presence Protocols}
\todo{Survey existing VR/AR co-presence protocols (OpenXR social layer, VRChat
world scripts, Roblox Luau, Horizon Worlds APIs) before making a comparative
claim about capability-scoped negotiation at thresholds.}

\subsection{CRDT-Based Shared State}
The Yjs~\cite{yjs2024} and Automerge~\cite{automerge2024}
libraries provide CRDT-based shared editing for documents.
Loro~\cite{loro2024} extends this model with a tree container
suited to hierarchical spatial state.  The \hm{}/HoloScript state layer uses
Loro as its state substrate without modification (\S\ref{sec:arch}).

\subsection{Agent Authorization Models}
The A2A agent card protocol~\cite{a2agooglespec2025} provides agent-card and task
interaction surfaces. The artifact includes a library-level adapter that
validates a card, derives proposed scopes from skill tags, and produces a
portal request-entry payload; production threshold orchestration remains unwired.

\subsection{Audit and Provenance}
Hash-chain provenance for agent-environment loops is formalized in
CAEL~\cite{cael2026}, which applies the chain to physics simulation steps. The
portal-entry component instead hashes receipt content and carries supplied
state-anchor fields. An analogous per-mutation world-state chain at co-presence
boundaries remains future work. The scope-based trust model relates to the
algebraic trust lattice of~\cite{paper29algebraictrust}.

% ── §3 System Architecture ───────────────────────────────────────────────────
\section{System Architecture}
\label{sec:arch}

\todo{SKELETON: full architecture section.  Current stubs below name the modules
shipped in HoloScript commits b111b6c2b–2b973eaaf.}

\begin{figure}[t]
  \centering
  \todo{Architecture diagram: show handshake/trait, MCP validator/store,
        optional WebRTC provider, and portal-entry receipt as separate components;
        use dashed edges for the unwired admission, transport, and receipt joins}
  \caption{HoloGate component overview. Solid connections denote separately
    implemented component interfaces; a future integration must bind the
    portal-presence admission record to the MCP intent path and choose a transport.
    The current MCP path uses in-process callback fan-out unless WebRTC is enabled
    explicitly. The portal-entry receipt carries an entry-state anchor, while
    mutation chaining is future work.}
  \label{fig:arch}
\end{figure}

\subsection{HoloPortal: Content-Negotiated Portal Protocol}
\label{sec:portal}

A \emph{portal} is represented as a typed spatial boundary with an admission
policy. The pure A2A library models card validation and scope derivation, while
the independent \texttt{@portalPresence} handler models entrant admission. A
bridge helper formats a request-entry payload, but no non-test caller currently
executes the two together. The library-level model has two conceptual steps:

\begin{enumerate}
\item \textbf{Card validation.} The helper validates a supplied A2A agent card
  and reads its self-described skill tags.
\item \textbf{Proposed scope derivation.} The helper derives a scope ladder
  from those tags and clamps it to caller-supplied portal maxima. A separate
  trait event later filters payload-supplied scopes against its own configuration.
\end{enumerate}

The separate \texttt{portal-entry.v1} schema can encode entrant,
representation, and granted-scope fields while extending the
\texttt{holotunnel.share-packet.v1} receipt schema
(commit~\texttt{b2da5af09}); no non-test caller currently builds that receipt
from a completed trait admission.

The \texttt{@portalPresence} handler (\texttt{af2b2a18e}) is exported from the
shared trait library and has targeted unit coverage. No non-test runtime
registration or consumer currently instantiates it, so general entity-level
availability remains an integration target.

\subsection{HoloDoor: Scope-Ladder Policy Enforcement}
\label{sec:scope}

\todo{Expand: allowedScopes gate, anchored wildcard matcher, reject/warn modes,
and the missing binding to portal admission and custody policy.}

The pure \texttt{validatePortalIntent} function validates a typed intent against
a supplied policy and requested scope without an external policy-service round
trip. Its MCP wrapper accepts both values from the tool caller; it does not read
the \texttt{@portalPresence} entrant registry. Consequently, this is a component
validator rather than enforcement of every admitted participant's intent.

For \texttt{mutate-zone}, scope matching is an anchored full-string match in
which \texttt{*} denotes any run of characters. Any configured glob match is
sufficient; there is no specificity ranking or winner selection. An intent that
fails policy is rejected by default or admitted with a warning when
\texttt{onScopeViolation} is \texttt{warn}.
These semantics assume a schema-conformant \texttt{requestedScope}. The MCP
handler casts this argument without checking enum membership at runtime; when
\texttt{allowedScopes} is absent, an unknown string can bypass some rank
comparisons. This is an open input-validation gap, not an admitted capability.

\subsection{Loro CRDT State Layer}
\label{sec:crdt}

The MCP state-authority component stores entities in a Loro map and persists the
document after committed changes (commit \texttt{af53937c1}, replacing the prior
deep-merge path). It broadcasts field-level diffs to an in-process subscriber
set. A separate opt-in \texttt{LoroWebRTCProvider} can synchronize the same
document (commit~\texttt{3e82a2eb5}), but no non-test portal or server-startup
caller currently enables it.

\subsection{Portal-Entry Receipt and Proposed Chain Model}
\label{sec:receipts}

The shipped portal-entry schema records entrant, representation, granted scopes,
tunnel/world/commit provenance, and a caller-supplied entry snapshot hash. Its
\texttt{rewindAvailable} field is a declaration; the current implementation does
not build an append-only mutation chain or reconstruct prior world state.

A future production chain would need validation that recomputes the builder's
content-derived identifier, per-mutation before/after state hashes,
authenticated roots, complete mutation payloads, concurrency semantics, and
replay that reconstructs and verifies the target state. The unit-level harness
in \S\ref{sec:receipt-eval} explores only a test-local subset of that model.

% ── §4 Implementation ────────────────────────────────────────────────────────
\section{Implementation}
\label{sec:impl}

The component prototypes are implemented in TypeScript (strict mode, ES2020).
The implementation spans the following modules and tracked commits:

\begin{center}
\begin{tabular}{lll}
\toprule
Module & Commit & Purpose \\
\midrule
\texttt{holo-portal-intent.ts} & \texttt{b111b6c2b} & Pure intent validator + delta conversion \\
\texttt{networking-tools.ts}   & \texttt{b111b6c2b} & MCP wrapper + callback fan-out \\
WebRTC peer transport          & \texttt{3e82a2eb5} & Opt-in Loro sync provider \\
Loro CRDT integration          & \texttt{af53937c1} & Replaces LWW deepMerge \\
A2A handshake                  & \texttt{e3876de57} & Pure card validation + scope derivation \\
Receipt extension              & \texttt{b2da5af09} & Standalone \texttt{portal-entry.v1} builder \\
\texttt{@portalPresence} trait & \texttt{af2b2a18e} & Shared trait library with test coverage \\
\bottomrule
\end{tabular}
\end{center}

Commit~\texttt{2b973eaaf} contributes the exhaustive scope-policy suite and a
separate test-local receipt-model harness. The former exercises shipped policy
validation; the latter does not implement a production chain or rewind path.
The historical threshold module introduced by commit~\texttt{73dcaad51} was
deleted by the immediately following commit~\texttt{1bab68e1e}; it is therefore
not counted as a shipped module. No current entry point wires all table rows
into a single admission-to-mutation pipeline.

\todo{SKELETON: add package placement note (packages/studio src/lib/ or
packages/core?), API surface overview, and HoloScript .hsplus policy snippet
showing a portal admission contract.}

% ── §5 Evaluation ────────────────────────────────────────────────────────────
\section{Evaluation}
\label{sec:eval}

All numeric timings below are from a historical standalone HoloGate-named
benchmark harness. The tracked receipt was added in commit
\texttt{9246bae47af19ffc8eed696221d8fa885416c90a}; its embedded
\texttt{commit} field identifies the benchmarked checkout as
\texttt{21da7eb86}. The receipt is \receiptid{hgbench\_mpgp6dt3} at
\texttt{scripts/receipts/hologate-portal-bench-receipt.json}. In this draft the
evidence is attributed to concrete
surfaces: an inline mirror of the intent validator, a pure threshold-state
simulation, a test-local callback set, and a standalone canonical-hash helper.
The harness did not invoke the integrated MCP/trait/WebRTC path. Hardware:
11th Gen Intel Core i7-11800H @ 2.30\,GHz (16 logical cores),
Node.js v24.15.0; from benchmark receipt \receiptid{hgbench\_mpgp6dt3}.
\measuredFrom{HoloScript/scripts/receipts/hologate-portal-bench-receipt.json}

\subsection{Micro-Benchmarks}
\label{sec:microbench}

\paragraph{Intent validation throughput.}
\measuredFrom{HoloScript/scripts/receipts/hologate-portal-bench-receipt.json}
The harness's inline mirror of \texttt{validatePortalIntent} processes
\textbf{772\,ns} average per validation
(\textbf{1.3\,M validations/sec}). This establishes a historical pure-function
component cost, not the cost or security of an admission-bound production path.

\paragraph{Portal threshold negotiation overhead.}
\measuredFrom{HoloScript/scripts/receipts/hologate-portal-bench-receipt.json}
The harness's pure in-process model of lane selection, scope resolution, receipt
header construction, and an arrival-event object takes
\textbf{1.6\,µs average} (p99: \textbf{4.0\,µs})
per iteration. It does not execute the current A2A handshake, portal-presence
trait, receipt builder, MCP tool, or any network crossing.

\paragraph{Delta broadcast fan-out.}
\measuredFrom{HoloScript/scripts/receipts/hologate-portal-bench-receipt.json}
Invoking a three-delta callback against a test-local \texttt{Set} of
\textbf{50 subscribers} costs
\textbf{2.2\,µs per broadcast event}.
This mirrors the in-process subscriber-registry algorithm; it does not commit a
Loro change or include WebRTC, WebSocket, serialization, or network delivery.

\paragraph{World-state SHA-256 cost.}
The standalone receipt helper benchmarks canonical-JSON SHA-256 over
\textbf{1\,000 synthetic entities}; hashing completes in
\textbf{2.0\,ms}.\measuredFrom{HoloScript/scripts/receipts/hologate-portal-bench-receipt.json}
This measures full-state hashing in the standalone harness, not a current
per-mutation receipt path. Incremental hashing and integration are future work.

\subsection{Loro CRDT Convergence}
\label{sec:crdt-eval}

One tracked unit test constructs two independent Loro documents, exchanges their
updates in-process, verifies identical merged state, and asserts that this single
small-entity round trip stays below a generous 10\,ms threshold.
\measuredFrom{HoloScript/packages/mcp-server/src/__tests__/holo-portal-intent.test.ts}
This is a regression threshold, not a latency distribution or a WebRTC network
benchmark; the earlier standalone CRDT performance receipt is absent.

\subsection{Multi-Entrant Conformance Harness}
\label{sec:scale}

The tracked harness constructs fixtures for 10 agent and 10 human entrants and
contains 32 hand-authored multi-component unit cases for lane selection, scope
enforcement, receipt validity and uniqueness, delta delivery, and concurrent
state updates. Lane and receipt records are built directly from helpers;
mutation and delivery use one in-process store and callback set, without
PortalPresence or WebRTC.
\measuredFrom{HoloScript/packages/mcp-server/src/__tests__/hologate-multi-entrant.test.ts}
It is conformance coverage, not a throughput or scalability measurement.

\subsection{Correctness: Scope Violation Rejection}
\label{sec:scope-eval}

The scope violation suite defines 32 hand-authored conformance cases:
12 scope$\times$intent cases under one fixed policy and 20 policy, glob,
warn-mode, and default-policy edge cases.
\measuredFrom{HoloScript/packages/mcp-server/src/__tests__/hologate-scope-rejection-exhaustive.test.ts}
\begin{itemize}
\item Full scope$\times$intent Cartesian product (\textit{all allowed/rejected combinations})
\item \texttt{allowedScopes} gate bypass attempts
\item Glob matching: prefix anchoring, trailing wildcard, deep paths, and grab \texttt{targetId} selection
\item Warn mode: correct flag-without-reject behavior
\end{itemize}
This is deterministic policy-conformance coverage, not an empirical
false-positive or false-negative rate.

\subsection{Correctness: In-Process Receipt-Model Harness}
\label{sec:receipt-eval}

The tracked unit-level harness contains eight passing in-process cases. Its
test-local model checks canonical SHA-256 determinism for sample world state,
hash change after an allowed mutation, an exit-to-entry identifier
back-reference, preservation of an entry-state hash and ordered mutation
metadata, rejection without mutation, insertion-order-independent hashing,
repeatable payload hashing, and separation of two agents' locally constructed
records.
\measuredFrom{HoloScript/packages/mcp-server/src/__tests__/hologate-audit-receipt-chain.test.ts}
The harness does not exercise a production append-only hash chain,
per-mutation before/after state hashes, authenticated tamper detection,
concurrent chain merging, or state reconstruction/rewind; those remain future
work.

% ── §6 Discussion ────────────────────────────────────────────────────────────
\section{Discussion}
\label{sec:discussion}

\subsection{Threat Model and Limitations}
\label{sec:threats}

\paragraph{Untrusted scope inputs.}
The handshake helper derives scopes from self-described card skill tags, while
the MCP path accepts caller-supplied scope and policy. Neither decision is
cryptographically bound to server-held admission state. The missing runtime
enum check on \texttt{requestedScope} further means the handler cannot rely on
its declarative tool schema as enforcement. Production integration must validate
the input, select policy server-side, and bind the resulting grant to an
authenticated entrant session.

\paragraph{Unwired enforcement boundary.}
The MCP \texttt{push\_portal\_intent} tool accepts \texttt{requestedScope} and
\texttt{spatialPolicy} from its caller; it neither authenticates those values
against the portal-presence entrant registry nor requires a prior admission.
Moreover, \texttt{push\_state\_delta} remains available under the same
\texttt{tools:write} authorization scope and bypasses the typed validator. The
current components therefore do not establish end-to-end admission-bound
enforcement. Production integration must derive policy and scope from trusted
server-held admission state and remove or separately authorize the raw path.

\paragraph{Production receipt boundary.}
The shipped portal receipt accepts a caller-supplied snapshot anchor and does not
provide an authenticated append-only mutation chain or replay implementation.
Production deployment must recompute and verify the builder's content-derived
identifier, bind before/after state hashes to each mutation, authenticate chain
roots, and verify state reconstruction before claiming tamper detection or rewind.

\paragraph{CRDT divergence under network partition.}
Loro's CRDT guarantees convergence only after all deltas are delivered.
Under indefinite partition, local state diverges.  The current model has no
partition-detection or explicit merge-conflict surfacing for human users.

\paragraph{Scale ceiling.}
The 20-entrant fixture is the largest tracked conformance harness, but it does
not establish throughput or a scalability ceiling. Future work must characterize
wall time, per-entrant latency, and fan-out behavior as $N$ grows.

\paragraph{Quest 3 validation.}
The evaluation is a Node.js component harness; it does not exercise a browser,
WebGPU, or a headset. Quest 3 validation is an open gate: the QuestProofPanel
(\texttt{packages/studio/src/components/quest/QuestProofPanel.tsx},
commit~\texttt{cbe4e77f3}) is ready but no device receipt has been recorded
in \texttt{.bench-logs/} at time of writing.

\todo{SKELETON: add subsection on applicability beyond spatial worlds
(document co-editing with agent participants, workflow co-execution).}

% ── §7 Conclusion ────────────────────────────────────────────────────────────
\section{Conclusion}
\label{sec:conclusion}

The repository contains the principal components for representing spatial
admission, validating typed intents, storing shared state in Loro, optionally
transporting Loro updates, and recording a portal-entry schema; \hg{} remains the
umbrella name for this design, not a measured integrated runtime surface. A
historical standalone harness measured an inline validator mirror at
\measuredFrom{HoloScript/scripts/receipts/hologate-portal-bench-receipt.json}772\,ns
and a threshold simulation at
\measuredFrom{HoloScript/scripts/receipts/hologate-portal-bench-receipt.json}1.6\,µs
average. Those values motivate integration experiments but do not establish
end-to-end latency or enforcement. The portal-entry receipt carries an
entry-state anchor; authenticated tamper detection and deterministic rewind
remain future work.

The immediate research result is therefore a tested component substrate and a
precise integration contract: bind server-held admission to every mutation
path, enable and measure a concrete transport, and derive verifiable receipt
state before claiming auditable co-presence.

\todo{SKELETON: Expand conclusion with 2--3 sentences on open questions
(Quest device validation, partition handling, cryptographic card binding)
and connection to the broader HoloScript provenance program.}

% ── Acknowledgements ─────────────────────────────────────────────────────────
% (omit for blind review)

% ── References ───────────────────────────────────────────────────────────────
\bibliographystyle{ACM-Reference-Format}
\bibliography{holoscript}
% Key bib entries needed (add to holoscript.bib before submission):
%   paper29algebraictrust   — Paper 29 Algebraic Trust (cite key TBD)
%   paper32crdtspatial      — Paper 32 CRDT spatial state (cite key TBD)
%   paper26latenttoken      — Paper 26 Latent Token Compression (cite key TBD)
%   cael2026                — Paper 0c CAEL (cite key: verify in holoscript.bib)
%   yjs2022                 — Yjs CRDT library
%   automerge2019           — Automerge CRDT library
%   loro2024                — Loro CRDT library
%   a2agooglespec2025       — Google A2A agent card protocol
%   eip712anchor2026        — EIP-712 chain anchor (if applicable)
% \todo{Fill all bib entries; run pnpm check:paper-claim-map before submission}

\end{document}
